
\documentclass{article}
\usepackage{colortbl}
\usepackage{makecell}
\usepackage{multirow}
\usepackage{supertabular}

\begin{document}

\newcounter{utterance}

\centering \large Interaction Transcript for game `dond', experiment `semi', episode 13 with Qwen3.6{-}35B{-}A3B{-}FP8{-}without{-}reasoning.
\vspace{24pt}

{ \footnotesize  \setcounter{utterance}{1}
\setlength{\tabcolsep}{0pt}
\begin{supertabular}{c@{$\;$}|p{.15\linewidth}@{}p{.15\linewidth}p{.15\linewidth}p{.15\linewidth}p{.15\linewidth}p{.15\linewidth}}
    \# & $\;$A & \multicolumn{4}{c}{Game Master} & $\;\:$B\\
    \hline

    \theutterance \stepcounter{utterance}  
    & & \multicolumn{4}{p{0.6\linewidth}}{
        \cellcolor[rgb]{0.9,0.9,0.9}{
            \makecell[{{p{\linewidth}}}]{
                \texttt{\tiny{[P1$\langle$GM]}}
                \texttt{Žaidžiate kooperatyvinį derybų žaidimą, kuriame turite sutarti su kitu žaidėju, kaip padalinti daiktų rinkinį tarp jūsų.} \\
\\ 
\texttt{Taisyklės:} \\
\texttt{(a) Jūs ir kitas žaidėjas gaunate daiktų rinkinį. Kiekvienam iš jūsų taip pat suteikiama slapta vertės funkcija, vaizduojanti, kiek vertinate kiekvieną daiktų tipą.} \\
\texttt{(b) Jūs keičiatės žinutėmis su kitu žaidėju, kad sutartumėte, kas gauna kokius daiktus. Kiekvienas gali išsiųsti ne daugiau nei 5 žinučių arba pabeigti žaidimą anksčiau, bet kuriuo metu pateikdamas slapią pasiūlymą.} \\
\texttt{(c) Jūs prašomi pateikti slaptą pasiūlymą, nurodantį norimus daiktus, suformatuotus skliaustuose taip: "{[}Pasiūlymas: <skaičius> <daikto pavadinimas>, <skaičius> <daikto pavadinimas>, <...>{]}"} \\
\texttt{(d) Jei jūsų pasiūlymai yra papildomi, t. y. yra pakankamai daiktų abiem pasiūlymams įvykdyti, kiekvienas žaidėjas gauna balus, remiantis gautų daiktų verčių suma. Priešingu atveju jūs abu gausite nulinį balą.} \\
\texttt{(e) Jūsų tikslas yra maksimizuoti savo balo ir kito žaidėjo balo sumą.} \\
\\ 
\texttt{Pradėkime.} \\
\\ 
\texttt{Prieinamų daiktų rinkinys yra:} \\
\\ 
\texttt{1 žiedas, 2 knygos, 2 popieriai, 1 bulvė} \\
\\ 
\texttt{Jūsų slapta vertės funkcija:} \\
\\ 
\texttt{žiedas: 5, knyga: 1, popierius: 0, bulvė: 3} \\
\\ 
\texttt{SVARBU: Jūsų žinutės, išskyrus slaptą pasiūlymą, perduodamos tiesiogiai kitam žaidėjui, todėl neįtraukite jokio atsakymo į taisykles ar teksto, kuriame nurodomas jūsų pranešimas. Norėdami pateikti slaptą pasiūlymą, naudokite nurodytą formatą. Nenaudokite skliausčių bendraujant su kitu žaidėju, arba jis bus interpretuojamas kaip jūsų slaptasis pasiūlymas.} \\
            }
        }
    }
    & \\ \\

    \theutterance \stepcounter{utterance}  
    & \multicolumn{4}{p{0.6\linewidth}}{
        \cellcolor[rgb]{0.9,0.9,0.9}{
            \makecell[{{p{\linewidth}}}]{
                \texttt{\tiny{[P1$\rangle$GM]}}
                \texttt{Sveiki! Supratau taisykles. Kadangi tai kooperatyvinis žaidimas, kurio tikslas yra maksimizuoti bendrą balų sumą, siūlau pradėti nuo atvirų derybų.} \\
\\ 
\texttt{Mano vertinimai:} \\
\texttt{{-} Žiedas: 5} \\
\texttt{{-} Knyga: 1} \\
\texttt{{-} Popierius: 0} \\
\texttt{{-} Bulvė: 3} \\
\\ 
\texttt{Jūsų vertinimai man kol kas nežinomi. Prašau pasidalinti savo vertinimais arba pasiūlyti savo norimus daiktus, kad galėtume suderinti optimalų padalijimą.} \\
\\ 
\texttt{Kokie yra jūsų vertinimai?} \\
            }
        }
    }
    & & \\ \\

    \theutterance \stepcounter{utterance}  
    & & & \multicolumn{2}{p{0.3\linewidth}}{
        \cellcolor[rgb]{0.9,0.9,0.9}{
            \makecell[{{p{\linewidth}}}]{
                \texttt{\tiny{[GM$|$GM]}}
                \texttt{continue} \\
            }
        }
    }
    & & \\ \\

    \theutterance \stepcounter{utterance}  
    & & & \multicolumn{2}{p{0.3\linewidth}}{
        \cellcolor[rgb]{0.9,0.9,0.9}{
            \makecell[{{p{\linewidth}}}]{
                \texttt{\tiny{[GM$|$GM]}}
                \texttt{Sveiki! Supratau taisykles. Kadangi tai kooperatyvinis žaidimas, kurio tikslas yra maksimizuoti bendrą balų sumą, siūlau pradėti nuo atvirų derybų.} \\
\\ 
\texttt{Mano vertinimai:} \\
\texttt{{-} Žiedas: 5} \\
\texttt{{-} Knyga: 1} \\
\texttt{{-} Popierius: 0} \\
\texttt{{-} Bulvė: 3} \\
\\ 
\texttt{Jūsų vertinimai man kol kas nežinomi. Prašau pasidalinti savo vertinimais arba pasiūlyti savo norimus daiktus, kad galėtume suderinti optimalų padalijimą.} \\
\\ 
\texttt{Kokie yra jūsų vertinimai?} \\
            }
        }
    }
    & & \\ \\

    \theutterance \stepcounter{utterance}  
    & & \multicolumn{4}{p{0.6\linewidth}}{
        \cellcolor[rgb]{0.9,0.9,0.9}{
            \makecell[{{p{\linewidth}}}]{
                \texttt{\tiny{[GM$\rangle$P2]}}
                \texttt{Žaidžiate kooperatyvinį derybų žaidimą, kuriame turite sutarti su kitu žaidėju, kaip padalinti daiktų rinkinį tarp jūsų.} \\
\\ 
\texttt{Taisyklės:} \\
\texttt{(a) Jūs ir kitas žaidėjas gaunate daiktų rinkinį. Kiekvienam iš jūsų taip pat suteikiama slapta vertės funkcija, vaizduojanti, kiek vertinate kiekvieną daiktų tipą.} \\
\texttt{(b) Jūs keičiatės žinutėmis su kitu žaidėju, kad sutartumėte, kas gauna kokius daiktus. Kiekvienas gali išsiųsti ne daugiau nei 5 žinučių arba pabeigti žaidimą anksčiau, bet kuriuo metu pateikdamas slapią pasiūlymą.} \\
\texttt{(c) Jūs prašomi pateikti slaptą pasiūlymą, nurodantį norimus daiktus, suformatuotus skliaustuose taip: "{[}Pasiūlymas: <skaičius> <daikto pavadinimas>, <skaičius> <daikto pavadinimas>, <...>{]}"} \\
\texttt{(d) Jei jūsų pasiūlymai yra papildomi, t. y. yra pakankamai daiktų abiem pasiūlymams įvykdyti, kiekvienas žaidėjas gauna balus, remiantis gautų daiktų verčių suma. Priešingu atveju jūs abu gausite nulinį balą.} \\
\texttt{(e) Jūsų tikslas yra maksimizuoti savo balo ir kito žaidėjo balo sumą.} \\
\\ 
\texttt{Pradėkime.} \\
\\ 
\texttt{Prieinamų daiktų rinkinys yra:} \\
\\ 
\texttt{1 žiedas, 2 knygos, 2 popieriai, 1 bulvė} \\
\\ 
\texttt{Jūsų slapta vertės funkcija:} \\
\\ 
\texttt{žiedas: 0, knyga: 3, popierius: 1, bulvė: 2} \\
\\ 
\texttt{SVARBU: Jūsų žinutės, išskyrus slaptą pasiūlymą, perduodamos tiesiogiai kitam žaidėjui, todėl neįtraukite jokio atsakymo į taisykles ar teksto, kuriame nurodomas jūsų pranešimas. Norėdami pateikti slaptą pasiūlymą, naudokite nurodytą formatą. Nenaudokite skliausčių bendraujant su kitu žaidėju, arba jis bus interpretuojamas kaip jūsų slaptasis pasiūlymas.} \\
\\ 
\texttt{Sveiki! Supratau taisykles. Kadangi tai kooperatyvinis žaidimas, kurio tikslas yra maksimizuoti bendrą balų sumą, siūlau pradėti nuo atvirų derybų.} \\
\\ 
\texttt{Mano vertinimai:} \\
\texttt{{-} Žiedas: 5} \\
\texttt{{-} Knyga: 1} \\
\texttt{{-} Popierius: 0} \\
\texttt{{-} Bulvė: 3} \\
\\ 
\texttt{Jūsų vertinimai man kol kas nežinomi. Prašau pasidalinti savo vertinimais arba pasiūlyti savo norimus daiktus, kad galėtume suderinti optimalų padalijimą.} \\
\\ 
\texttt{Kokie yra jūsų vertinimai?} \\
            }
        }
    }
    & \\ \\

    \theutterance \stepcounter{utterance}  
    & & & \multicolumn{4}{p{0.6\linewidth}}{
        \cellcolor[rgb]{0.9,0.9,0.9}{
            \makecell[{{p{\linewidth}}}]{
                \texttt{\tiny{[GM$\langle$P2]}}
                \texttt{Sveiki! Ačiū už pasiūlymą bendradarbiauti. Kadangi tikslas yra maksimizuoti bendrą naudą, siūlau paskirstyti daiktus taip, kad kiekvienas gautų tai, kas jam vertingiausia.} \\
\\ 
\texttt{Remiantis jūsų pateiktais vertinimais (žiedas: 5, bulvė: 3) ir mano slaptais vertinimais (knyga: 3, bulvė: 2, popierius: 1), optimaliausias variantas būtų:} \\
\\ 
\texttt{1. **Jūs gaunate:** 1 žiedą (jums vertinga 5) ir 1 bulvę (jums vertinga 3). Jūsų bendras pelnas: 8.} \\
\texttt{2. **Aš gaunu:** 2 knygas (man vertinga 3 kiekviena, iš viso 6) ir 2 popierius (man vertinga 1 kiekvienas, iš viso 2). Mano bendras pelnas: 8.} \\
\\ 
\texttt{Bendra suma: 16 balų.} \\
\\ 
\texttt{Patikrinkite, ar šis padalijimas jums tinka. Jei taip, prašau pateikti savo **slaptą pasiūlymą** nurodytu formatu, kad galėtume užbaigti žaidimą.} \\
\\ 
\texttt{Jei norite pakeisti dalį, galime aptarti. Pavyzdžiui, jei man atiduosite vieną knygą, o jūs man duosite bulvę, bendra suma sumažėtų, nes man bulvė vertinga mažiau nei knyga, o jums žiedas yra labai vertingas.} \\
\\ 
\texttt{Laukiu jūsų slaptą pasiūlymo, jei sutinkate su šiuo planu.} \\
            }
        }
    }
     \\ \\

    \theutterance \stepcounter{utterance}  
    & & & \multicolumn{2}{p{0.3\linewidth}}{
        \cellcolor[rgb]{0.9,0.9,0.9}{
            \makecell[{{p{\linewidth}}}]{
                \texttt{\tiny{[GM$|$GM]}}
                \texttt{continue} \\
            }
        }
    }
    & & \\ \\

    \theutterance \stepcounter{utterance}  
    & & & \multicolumn{2}{p{0.3\linewidth}}{
        \cellcolor[rgb]{0.9,0.9,0.9}{
            \makecell[{{p{\linewidth}}}]{
                \texttt{\tiny{[GM$|$GM]}}
                \texttt{Sveiki! Ačiū už pasiūlymą bendradarbiauti. Kadangi tikslas yra maksimizuoti bendrą naudą, siūlau paskirstyti daiktus taip, kad kiekvienas gautų tai, kas jam vertingiausia.} \\
\\ 
\texttt{Remiantis jūsų pateiktais vertinimais (žiedas: 5, bulvė: 3) ir mano slaptais vertinimais (knyga: 3, bulvė: 2, popierius: 1), optimaliausias variantas būtų:} \\
\\ 
\texttt{1. **Jūs gaunate:** 1 žiedą (jums vertinga 5) ir 1 bulvę (jums vertinga 3). Jūsų bendras pelnas: 8.} \\
\texttt{2. **Aš gaunu:** 2 knygas (man vertinga 3 kiekviena, iš viso 6) ir 2 popierius (man vertinga 1 kiekvienas, iš viso 2). Mano bendras pelnas: 8.} \\
\\ 
\texttt{Bendra suma: 16 balų.} \\
\\ 
\texttt{Patikrinkite, ar šis padalijimas jums tinka. Jei taip, prašau pateikti savo **slaptą pasiūlymą** nurodytu formatu, kad galėtume užbaigti žaidimą.} \\
\\ 
\texttt{Jei norite pakeisti dalį, galime aptarti. Pavyzdžiui, jei man atiduosite vieną knygą, o jūs man duosite bulvę, bendra suma sumažėtų, nes man bulvė vertinga mažiau nei knyga, o jums žiedas yra labai vertingas.} \\
\\ 
\texttt{Laukiu jūsų slaptą pasiūlymo, jei sutinkate su šiuo planu.} \\
            }
        }
    }
    & & \\ \\

    \theutterance \stepcounter{utterance}  
    & & \multicolumn{4}{p{0.6\linewidth}}{
        \cellcolor[rgb]{0.9,0.9,0.9}{
            \makecell[{{p{\linewidth}}}]{
                \texttt{\tiny{[P1$\langle$GM]}}
                \texttt{Sveiki! Ačiū už pasiūlymą bendradarbiauti. Kadangi tikslas yra maksimizuoti bendrą naudą, siūlau paskirstyti daiktus taip, kad kiekvienas gautų tai, kas jam vertingiausia.} \\
\\ 
\texttt{Remiantis jūsų pateiktais vertinimais (žiedas: 5, bulvė: 3) ir mano slaptais vertinimais (knyga: 3, bulvė: 2, popierius: 1), optimaliausias variantas būtų:} \\
\\ 
\texttt{1. **Jūs gaunate:** 1 žiedą (jums vertinga 5) ir 1 bulvę (jums vertinga 3). Jūsų bendras pelnas: 8.} \\
\texttt{2. **Aš gaunu:** 2 knygas (man vertinga 3 kiekviena, iš viso 6) ir 2 popierius (man vertinga 1 kiekvienas, iš viso 2). Mano bendras pelnas: 8.} \\
\\ 
\texttt{Bendra suma: 16 balų.} \\
\\ 
\texttt{Patikrinkite, ar šis padalijimas jums tinka. Jei taip, prašau pateikti savo **slaptą pasiūlymą** nurodytu formatu, kad galėtume užbaigti žaidimą.} \\
\\ 
\texttt{Jei norite pakeisti dalį, galime aptarti. Pavyzdžiui, jei man atiduosite vieną knygą, o jūs man duosite bulvę, bendra suma sumažėtų, nes man bulvė vertinga mažiau nei knyga, o jums žiedas yra labai vertingas.} \\
\\ 
\texttt{Laukiu jūsų slaptą pasiūlymo, jei sutinkate su šiuo planu.} \\
            }
        }
    }
    & \\ \\

    \theutterance \stepcounter{utterance}  
    & \multicolumn{4}{p{0.6\linewidth}}{
        \cellcolor[rgb]{0.9,0.9,0.9}{
            \makecell[{{p{\linewidth}}}]{
                \texttt{\tiny{[P1$\rangle$GM]}}
                \texttt{Šis pasiūlymas man tinka. Jis leidžia mums abiem gauti maksimalią naudą iš turimų resursų.} \\
\\ 
\texttt{Pateikiu savo slaptą pasiūlymą:} \\
\\ 
\texttt{{[}Pasiūlymas: 1 žiedas, 1 bulvė{]}} \\
            }
        }
    }
    & & \\ \\

    \theutterance \stepcounter{utterance}  
    & & & \multicolumn{2}{p{0.3\linewidth}}{
        \cellcolor[rgb]{0.9,0.9,0.9}{
            \makecell[{{p{\linewidth}}}]{
                \texttt{\tiny{[GM$|$GM]}}
                \texttt{proposal} \\
            }
        }
    }
    & & \\ \\

    \theutterance \stepcounter{utterance}  
    & & & \multicolumn{2}{p{0.3\linewidth}}{
        \cellcolor[rgb]{0.9,0.9,0.9}{
            \makecell[{{p{\linewidth}}}]{
                \texttt{\tiny{[GM$|$GM]}}
                [1, 0, 0, 1]
            }
        }
    }
    & & \\ \\

    \theutterance \stepcounter{utterance}  
    & & \multicolumn{4}{p{0.6\linewidth}}{
        \cellcolor[rgb]{0.9,0.9,0.9}{
            \makecell[{{p{\linewidth}}}]{
                \texttt{\tiny{[GM$\rangle$P2]}}
                \texttt{Kitas žaidėjas pateikė slaptą pasiūlymą. Dabar jūsų eilė pateikti savo pasiūlymą.} \\
            }
        }
    }
    & \\ \\

    \theutterance \stepcounter{utterance}  
    & & & \multicolumn{4}{p{0.6\linewidth}}{
        \cellcolor[rgb]{0.9,0.9,0.9}{
            \makecell[{{p{\linewidth}}}]{
                \texttt{\tiny{[GM$\langle$P2]}}
                \texttt{{[}Slaptas pasiūlymas: 2 knygos, 2 popieriai{]}} \\
            }
        }
    }
     \\ \\

    \theutterance \stepcounter{utterance}  
    & & & \multicolumn{2}{p{0.3\linewidth}}{
        \cellcolor[rgb]{0.9,0.9,0.9}{
            \makecell[{{p{\linewidth}}}]{
                \texttt{\tiny{[GM$|$GM]}}
                \texttt{proposal does not start with "pasiūlymas:"} \\
            }
        }
    }
    & & \\ \\

\end{supertabular}
}

\end{document}
